% ------------------------------------------------------------------
%  Capitolo 13 — Insegnare per imparare
%  Parte III
%  Ricerca di riferimento: ricerca 01 §5.8, §6.3 (studio in gruppo);
%    03 §3.4, §6
%  Voci bibliografiche principali: nestojko2014, fiorella2013, koh2018,
%    rozenblit2002
%  Epigrafe: ✔ verificata sul testo (vedi ricerca 03 e registro, ricerca 00)
%  Scheda del capitolo: ricerca/06_indice-ragionato.md, capitolo 13
%
%  STATO: prima stesura (23/09/2026), circa 2.200 parole.
%  Dati verificati: S-14 (Koh et al. 2018: 124 studenti, quattro
%  condizioni, test a una settimana); B-08 (fiorella2013). Fiorella &
%  Mayer 2013 verificato sull'abstract il 23/09/2026 (a una settimana
%  solo il gruppo che aveva insegnato davvero supera il controllo).
%  Nestojko et al. 2014: risultato qualitativo (più ricordo, soprattutto
%  delle idee principali, e migliore organizzazione).
%  Il tutoraggio tra pari (scheda: «tutoraggio tra pari») è rinviato ai
%  capp. 20 e 23: manca ancora in bibliografia una fonte specifica.
%  La scena d'apertura è tipica, non autobiografica (decisione 10).
% ------------------------------------------------------------------
\chapter{Insegnare per imparare}
\label{cap:insegnare}

\epigrafe[gli uomini, mentre insegnano, imparano]{homines dum docent discunt}{Seneca, Epistulae ad Lucilium 7, 8}

\iniziale{M}{artedì sera}. Chiara, studentessa del secondo anno di
biotecnologie, ha appena ricevuto un messaggio da Marco, il compagno con
cui di solito segue le lezioni: è stato a casa con la febbre, ha perso
la lezione di genetica sul controllo dell'espressione genica e le chiede
se domani, prima della lezione, può spiegargliela in dieci minuti. Chiara
risponde di sì senza pensarci: era presente, ha preso appunti, le è
sembrato tutto chiaro.

Poi, dopo cena, prova a immaginare come comincerà la spiegazione. E si
accorge che non è così semplice. Sa che cos'è un operone, ma non sa
bene come dirlo in due frasi; ricorda che il lattosio \enquote{accende}
qualcosa, ma non ricorda più se il repressore si stacca o si attacca;
e quando Marco le chiederà -- perché glielo chiederà -- \enquote{sì, ma
perché la cellula fa così?}, non è sicura di saper rispondere. Riapre
gli appunti, questa volta con domande precise in testa. Rilegge il
paragrafo del manuale che aveva saltato. Si fa uno schema. Prova a
spiegare tutto ad alta voce, da sola, in camera, e alla seconda volta ci
riesce. La mattina dopo la spiegazione a Marco va bene. Ma chi ha
imparato di più, in questa storia, è probabilmente Chiara.

Seneca lo aveva osservato quasi duemila anni fa. Scrivendo a Lucilio a
proposito delle compagnie da frequentare, gli consiglia di stare con chi
può renderlo migliore e di accogliere chi lui stesso può rendere
migliore, perché \enquote{queste cose avvengono reciprocamente: gli
uomini, mentre insegnano, imparano}.\footnote{Seneca, \emph{Epistulae
ad Lucilium} 7, 8: \emph{mutuo ista fiunt, et homines dum docent
discunt}.} Questo capitolo si chiede perché insegnare faccia imparare
chi insegna, e come sfruttare questo effetto quando si studia, da soli
o in gruppo.

\section{Docendo discimus}

Il motto latino che quasi tutti conoscono, \emph{docendo discimus},
\enquote{insegnando impariamo}, non si trova in nessun autore antico:
è una formula moderna, ricavata proprio dalla frase di
Seneca.\footnote{Non è l'unico caso: anche \emph{repetita iuvant}, che
molti credono un detto classico, non si trova negli autori antichi; la
sua origine lontana è un verso di Orazio (\emph{Ars poetica} 365) che
parla di opere d'arte che piacciono di più a ogni rilettura, non di
memoria.} Ma l'idea è antica quanto l'insegnamento. Chiunque abbia
dovuto spiegare qualcosa -- un docente alla prima lezione di un corso
nuovo, un dottorando alla sua prima esercitazione, un fratello maggiore
alle prese con i compiti di matematica del più piccolo -- conosce la
sensazione di aver capito davvero un argomento solo nel momento in cui
ha dovuto farlo capire a un altro.

Perché accade? Una prima ragione riguarda un'illusione che tutti
condividiamo. Due psicologi di Yale, Leonid Rozenblit e Frank Keil,
chiesero ad alcune persone di valutare quanto bene capissero il
funzionamento di oggetti comuni, come una cerniera lampo o lo sciacquone
del bagno. Poi chiesero loro di spiegarne il funzionamento per iscritto,
passo per passo. Dopo aver provato a spiegare, le stesse persone
abbassavano nettamente la valutazione della propria comprensione
\parencite{rozenblit2002}. I ricercatori chiamarono questo fenomeno
\emph{illusione della profondità esplicativa}: crediamo di capire le
cose molto meglio di quanto le capiamo, e ce ne accorgiamo soltanto
quando dobbiamo spiegarle. Chiara credeva di aver capito la lezione. Lo
ha scoperto non vero quando ha dovuto prepararsi a spiegarla.

Insegnare, dunque, è prima di tutto un modo di scoprire ciò che non
sai: un antidoto alle illusioni di competenza di cui abbiamo parlato nel
capitolo~\ref{cap:inganni}. Ma non è l'unica ragione per cui funziona,
e forse nemmeno la principale.

\section{Che cosa dicono le prove}

La ricerca sperimentale ha scomposto l'insegnare nelle sue parti, e ha
cercato di capire quale di esse faccia imparare.

La prima parte è l'\emph{aspettativa} di dover insegnare. In uno studio
di John Nestojko e colleghi, alcuni studenti leggevano un testo sapendo
che dopo avrebbero dovuto sostenere un test; altri lo leggevano
credendo che avrebbero dovuto insegnarlo a qualcun altro
\parencite{nestojko2014}. In realtà nessuno insegnava: tutti facevano lo
stesso test. Ma chi si aspettava di insegnare ricordava di più, in
particolare le idee principali, e ricordava in modo più ordinato. Il
solo pensiero di dover spiegare cambia il modo in cui si legge: si
cerca la struttura, si distingue l'essenziale dal dettaglio, ci si
chiede come mettere in fila le cose.

La seconda parte è l'insegnare \emph{davvero}. Logan Fiorella e Richard
Mayer hanno confrontato tre gruppi di studenti: uno studiava un testo
come per un normale esame; un altro lo studiava per prepararsi a
insegnarlo; un terzo si preparava e poi lo spiegava davvero, registrando
una breve lezione davanti a una videocamera \parencite{fiorella2013}.
Subito dopo, entrambi i gruppi che si erano preparati a insegnare
andavano meglio del primo. Una settimana dopo, però, solo il gruppo che
aveva effettivamente spiegato conservava un vantaggio netto. Prepararsi
a insegnare aiuta; insegnare aiuta di più, e più a lungo.

Ma che cosa, esattamente, dell'insegnare fa imparare? Aloysius Koh,
Sze Chi Lee e Stephen Lim hanno risposto con un esperimento molto
pulito \parencite{koh2018}. Hanno diviso 124 studenti universitari in
quattro gruppi. Il primo insegnava leggendo da un testo preparato, come
un relatore che legge le slide; il secondo insegnava a memoria, senza
appunti; il terzo non insegnava affatto, ma faceva pratica del recupero,
scrivendo a libro chiuso ciò che ricordava; il quarto, di controllo,
passava lo stesso tempo facendo esercizi di aritmetica. Una settimana
dopo, il gruppo che aveva insegnato leggendo non andava meglio del
controllo. I due gruppi che andavano meglio erano quello che aveva
insegnato a memoria e quello che aveva soltanto fatto pratica del
recupero.

La conclusione è importante, e un po' sorprendente: il beneficio
dell'insegnare deriva in larga parte dal fatto che insegnare \emph{senza
appunti} obbliga a recuperare. È la pratica del recupero del
capitolo~\ref{cap:recupero}, travestita da lezione. A essa si
aggiungono l'organizzazione del materiale, favorita dall'aspettativa di
spiegarlo, e la scoperta delle lacune. Ma senza il recupero, l'effetto
quasi scompare.

\section{Spiegare senza appunti}

Da questi risultati viene un consiglio preciso: se vuoi imparare
insegnando, \emph{insegna a memoria}. Leggere a un compagno i propri
appunti, o ripetere ad alta voce le slide del docente con le slide
davanti, è una rilettura fatta a voce alta: rassicura chi parla,
annoia chi ascolta, e lascia poco a entrambi. Spiegare senza guardare,
invece, è faticoso, si inceppa, costringe a fermarsi, a tornare
indietro, a cercare le parole. Quella fatica è il segno che il recupero
sta avvenendo.

In pratica, la sequenza che funziona è questa. Prima studi
l'argomento, sapendo che dovrai spiegarlo: cerchi la struttura, le idee
principali, gli esempi migliori. Poi chiudi il libro e spieghi, ad alta
voce, a qualcuno o anche a nessuno: a un compagno, a un familiare
paziente, alla videocamera del telefono, alla sedia vuota davanti a te.
Infine riapri il materiale e controlli: che cosa hai saltato? dove ti
sei inceppato? che cosa hai detto di sbagliato? Quei punti tornano nel
tuo mazzo di domande (capitolo~\ref{cap:tempo}).

Un ascoltatore vero, se c'è, aggiunge qualcosa che la sedia vuota non
può dare: le domande. \enquote{Perché?}, \enquote{E se invece\dots?},
\enquote{Non ho capito il passaggio tra questo e quello}. Ogni domanda
obbliga chi spiega a elaborare, a collegare, a trovare un esempio: è
l'interrogazione elaborativa del capitolo~\ref{cap:capire}, fatta da
un altro al posto tuo. Per questo spiegare a qualcuno che ne sa meno di
te, ma che è curioso e non si accontenta, è uno degli esercizi di studio
più efficaci che esistano.

Per chi prepara un esame orale, tutto questo ha una traduzione
immediata. L'esame orale \emph{è} una spiegazione senza appunti,
davanti a qualcuno che fa domande. Allenarsi a spiegare ad alta voce,
a memoria, ogni argomento del programma, possibilmente davanti a un
compagno che fa la parte del docente, è l'allenamento più vicino alla
prova vera che si possa immaginare. Registrarsi e riascoltarsi, per
quanto imbarazzante, rivela le esitazioni, i giri di parole, i punti
in cui si parla molto per dire poco.

\section{Studiare in gruppo, bene}

Tutto ciò porta naturalmente allo studio in gruppo, una delle
abitudini più diffuse tra gli studenti universitari italiani, e una
delle più ambivalenti. Studiare con altri può essere molto efficace o
del tutto inutile, a seconda di ciò che si fa. La ricerca non dice che
il gruppo funzioni o non funzioni in sé; dice che funziona quando i
componenti si interrogano e si spiegano le cose a vicenda, e non
funziona quando si limitano a rileggere insieme, ad ascoltare chi ne sa
di più, o a chiacchierare.

Le forme inefficaci sono facili da riconoscere. C'è il gruppo che
\enquote{fa il riassunto}: uno legge ad alta voce il manuale, gli altri
seguono e sottolineano. C'è il gruppo del compagno bravo: uno spiega
tutto, gli altri ascoltano e annuiscono, e a imparare -- come ormai sai
-- è soprattutto quello che spiega. C'è il gruppo che si rassicura:
ci si racconta quanto si è indietro, si confrontano le ansie, si
commentano i docenti, e la serata passa. Nessuna di queste è pratica
del recupero.

Le forme efficaci hanno invece alcune caratteristiche comuni. Ciascuno
arriva avendo \emph{già studiato} la parte assegnata: il gruppo non
serve a studiare per la prima volta, ma a recuperare, verificare e
approfondire. I libri e gli appunti restano \emph{chiusi} durante le
spiegazioni e le interrogazioni, e si aprono solo per controllare. I
\emph{ruoli ruotano}: chi spiega, chi fa domande, chi controlla sul
testo; nessuno resta per tutta la sessione nel ruolo comodo
dell'ascoltatore. E il gruppo è \emph{piccolo}: in due, tre o al
massimo quattro, ciascuno parla abbastanza da trarne vantaggio.

\begin{inpratica}
\textbf{Una sessione di studio a coppie, in trenta minuti.} Prima della
sessione, ciascuno studia lo stesso argomento e prepara cinque domande,
alcune di memoria e alcune di applicazione.
\begin{enumerate}
  \item \emph{Primi dieci minuti}: A spiega a B, a libro chiuso, la
        prima metà dell'argomento; B ascolta, interrompe con domande
        (\enquote{perché?}, \enquote{fammi un esempio}) e alla fine
        controlla sul testo che cosa manca.
  \item \emph{Secondi dieci minuti}: si invertono i ruoli sulla seconda
        metà.
  \item \emph{Ultimi dieci minuti}: ciascuno pone all'altro le proprie
        cinque domande, a libro chiuso; le risposte si controllano
        insieme.
\end{enumerate}
Alla fine, ciascuno annota gli errori e le lacune da riprendere nella
sessione successiva, tra qualche giorno.
\end{inpratica}

Un'ultima considerazione, che va oltre il metodo. Lo studio in gruppo
fatto bene non è soltanto più efficace: è anche più umano. L'università
può essere un luogo molto solitario, soprattutto nei primi anni e
soprattutto per chi viene da lontano (capitolo~\ref{cap:transizione}).
Spiegare le cose agli altri, farsi spiegare le cose dagli altri,
scoprire insieme ciò che non si sapeva, è anche un modo di appartenere
a una comunità di persone che studiano. Seneca, nella lettera che apre
questo capitolo, parlava proprio di questo: delle compagnie che ci
rendono migliori e di quelle che noi possiamo rendere migliori. Si
impara di più quando si impara con altri e per altri.

C'è anche un'occasione che molti studenti lasciano passare. Diversi
atenei affidano attività di tutorato a studenti degli anni successivi,
che seguono le matricole negli esercizi o nel metodo di studio. Se ti
viene proposto, accetta: oltre ad aiutare chi comincia, è uno dei modi
migliori per consolidare ciò che hai studiato negli anni precedenti, e
per scoprire quanto ne ricordi davvero.

\begin{riquadro}{Per chi insegna}
Se spiegare fa imparare chi spiega, in molte lezioni universitarie a
imparare di più è il docente. Bastano pochi minuti per restituire una
parte di questo vantaggio agli studenti: fermare la lezione e chiedere
a ciascuno di spiegare al vicino, senza appunti, il passaggio appena
visto; porre una domanda concettuale, far rispondere tutti
individualmente, poi far discutere le risposte a coppie prima di
rivelare quella corretta. Se ne parla nel capitolo~\ref{cap:docenti}.
\end{riquadro}

\asterismo

Chiara, la sera prima di spiegare la lezione a Marco, ha fatto senza
saperlo tutto ciò che la ricerca consiglia: si è accorta di non sapere
ciò che credeva di sapere, ha studiato cercando la struttura, ha provato
a spiegare a memoria, ha controllato e ha riprovato. Non aveva bisogno
di un compagno ammalato per farlo. Chiunque studi può immaginare di
dover spiegare domani ciò che studia oggi, e poi spiegarlo davvero, a
voce alta, senza guardare. Insegnare, in fondo, è recuperare ad alta
voce.

\begin{daricordare}
Insegnare è recuperare ad alta voce. Per imparare insegnando, studia
come se dovessi spiegare, poi spiega a memoria, senza appunti, e alla
fine controlla. Nello studio in gruppo, interrogatevi e spiegatevi a
vicenda: non rileggete insieme.
\end{daricordare}

\begin{domande}
  \item Da dove viene il motto \emph{docendo discimus}?
  \item Che cos'è l'illusione della profondità esplicativa, e che
        cosa c'entra con lo studio?
  \item Descrivi l'esperimento di Koh e colleghi. Perché insegnare
        leggendo gli appunti rende poco?
  \item Quali caratteristiche distinguono uno studio di gruppo
        efficace da uno inutile?
  \item (Ripresa, capitolo~\ref{cap:immagini}) Quando conviene il
        metodo dei luoghi, e quando no?
  \item (Ripresa, capitolo~\ref{cap:recupero}) Quali sono i tre passi
        del ciclo di base della pratica del recupero?
\end{domande}

\begin{sintesi}
  \item Spiegare ad altri fa imparare chi spiega: lo osservava già
        Seneca, e la ricerca lo conferma.
  \item Aspettarsi di insegnare migliora l'organizzazione di ciò che si
        studia; insegnare davvero produce un vantaggio più duraturo.
  \item Il beneficio viene soprattutto dal recupero: insegnare a
        memoria funziona, insegnare leggendo gli appunti quasi no.
  \item Provare a spiegare rivela le lacune che la rilettura nasconde
        (illusione della profondità esplicativa).
  \item Lo studio in gruppo funziona se è piccolo, preparato, a libro
        chiuso e con ruoli che ruotano; non funziona se si rilegge
        insieme o si ascolta soltanto.
\end{sintesi}

\finecapitolo
